% Capítulo 5: Plan de pruebas, implementación y evaluación.
\chapter{Plan de pruebas, implementación y evaluación del sistema}
\label{cap:pruebas-implementacion-evaluacion}

El anteproyecto plantea una evaluación progresiva del prototipo, desde la
integración de sus componentes hasta los ensayos en el invernadero. El plan
comprende la implementación física, el despliegue de los modelos, las pruebas
controladas y el registro de evidencias para analizar la confiabilidad del sistema
y la viabilidad del procesamiento local.

En laboratorio se prevé comprobar la interoperabilidad de los componentes y el
consumo de recursos. Posteriormente se realizarán pruebas de campo para examinar
la detección directa de la plaga bajo las condiciones de iluminación, follaje y
operación del sitio de estudio. Los ajustes técnicos deberán registrarse junto
con la versión del prototipo y las condiciones de cada ensayo.

Antes de ejecutar las pruebas se deberán definir la unidad de evaluación, el
procedimiento de identificación experta de la plaga y el significado operativo
de detección temprana. También se establecerán los criterios de aceptación, las
métricas de detección y los indicadores de latencia, consumo de recursos y
disponibilidad de datos. Estos elementos permitirán separar el desempeño de la
percepción visual del aporte del contexto ambiental y del conocimiento experto.

\textit{[Pendiente: completar el protocolo, los casos de prueba y los criterios
de aceptación; después, incorporar la implementación y los registros de las
pruebas efectivamente ejecutadas. No se reportan resultados experimentales en
esta versión.]}
